\ifdefined\gkver\else\def\gkver{student}\fi
\documentclass[\gkver]{gaokaozhenti}
\begin{document}

\chapter{2018年北京理综}
%% paperType: 理综物理部分

\begin{enumerate}
\item
%% number: 13
%% paperName: 2018年北京理综第 13 题 6 分
%% typeId: 1
%% type: 单选题
%% chapter: 原子和原子核
%% point: 人工核反应
%% method: 无
%% score: 6
%% degree: 950
%% duplicateId: 0
%% body:
在核反应方程 \ce{^{4}_{2}He + ^{14}_{7}N -> ^{17}_{8}O + X} 中，X 表示的是\xzanswer{A}

\fourchoices[answer=A]
{质子}
{中子}
{电子}
{$\alpha$ 粒子}

%% answer: A
%% memo:
\memoanswer{
由电荷数守恒和质量数守恒得出 X 的电荷数和质量数，从而确定 X 的种类。设 X 为 $\ce{^{A}_{Z}X}$，由质量数守恒得 $4 + 14 = 17 + A$，则 $A = 1$；由电荷数守恒得 $2 + 7 = 8 + Z$，则 $Z = 1$，即 X 为 $\ce{^{1}_{1}H}$，为质子，故 A 正确。
}

\item
%% number: 14
%% paperName: 2018年北京理综第 14 题 6 分
%% typeId: 1
%% type: 单选题
%% chapter: 气体性质
%% point: 分子动理论
%% method: 无
%% score: 6
%% degree: 850
%% duplicateId: 0
%% body:
关于分子动理论，下列说法正确的是\xzanswer{C}

\fourchoices[answer=C]
{气体扩散的快慢与温度无关}
{布朗运动是液体分子的无规则运动}
{分子间同时存在着引力和斥力}
{分子间的引力总是随分子间距增大而增大}

%% answer: C
%% memo:
\memoanswer{
温度越高，扩散越快，故 A 错误；布朗运动反映了分子的无规则运动，但不是分子运动，故 B 错误；分子间同时存在引力和斥力，故 C 正确；分子间的引力总是随分子间距增大而减小，故 D 错误。
}

\item
%% number: 15
%% paperName: 2018年北京理综第 15 题 6 分
%% typeId: 1
%% type: 单选题
%% chapter: 光学
%% point: 光的干涉
%% method: 无
%% score: 6
%% degree: 850
%% duplicateId: 0
%% body:
用双缝干涉实验装置得到白光的干涉条纹，在光源与单缝之间加上红色滤光片后\xzanswer{D}

\fourchoices[answer=D]
{干涉条纹消失}
{彩色条纹中的红色条纹消失}
{中央条纹变成暗条纹}
{中央条纹变成红色}

%% answer: D
%% memo:
\memoanswer{
当用白光做干涉实验时，光屏上得到干涉条纹，中央部分是白色的亮条纹；在光源与单缝之间加上红色滤光片后，只有红光能通过单缝，在光屏上形成红色的干涉条纹，中央部分为红色的亮条纹，故 ABC 错误 D 正确。
}

\item
%% number: 16
%% paperName: 2018年北京理综第 16 题 6 分
%% typeId: 1
%% type: 单选题
%% chapter: 机械波
%% point: 波的多解性
%% method: 无
%% score: 6
%% degree: 750
%% duplicateId: 0
%% body:
如图所示，一列简谐横波向右传播，P、Q 两质点平衡位置相距 $0.15\Um$。当 P 运动到上方最大位移处时，Q 刚好运动到下方最大位移处，则这列波的波长可能是\xzanswer{B}

\onepicture[w=4.5cm]{\includesvg{figs/16}}

\fourchoices[answer=B]
{$0.60\Um$}
{$0.30\Um$}
{$0.20\Um$}
{$0.15\Um$}

%% answer: B
%% memo:
\memoanswer{
由题意可知 P、Q 两质点平衡位置与波长之间满足：$x_{PQ} = (n + 0.5)\lambda$，其中 $n = 0$、1、2……，得 $\lambda = \frac{0.3\Um}{2n + 1}$，只有 B 项合乎要求。
}

\item
%% number: 17
%% paperName: 2018年北京理综第 17 题 6 分
%% typeId: 1
%% type: 单选题
%% chapter: 曲线运动
%% point: 天体运动
%% method: 无
%% score: 6
%% degree: 650
%% duplicateId: 0
%% body:
若想检验“使月球绕地球运动的力”与“使苹果落地的力”遵循同样的规律，在已知月地距离约为地球半径 60 倍的情况下，需要验证\xzanswer{B}

\fourchoices[answer=B]
{地球吸引月球的力约为地球吸引苹果的力的 $1/60^{2}$}
{月球公转的加速度约为苹果落向地面加速度的 $1/60^{2}$}
{自由落体在月球表面的加速度约为地球表面的 $1/6$}
{苹果在月球表面受到的引力约为在地球表面的 $1/60$}

%% answer: B
%% memo:
\memoanswer{
“使月球绕地球运动的力”与“使苹果落地的力”遵循同样的规律即万有引力定律。已知月地距离约为地球半径 60 倍，由 $\frac{GMm}{r^{2}} = ma$，得 $a = \frac{GM}{r^{2}}$，所以月球围绕地球运动的加速度是地球表面重力加速度的 $\frac{1}{60^{2}}$，B 项正确；由于月球的质量与苹果的质量不等，所以由 $F = \frac{GMm}{r^{2}}$ 可知 A 错误；由题目的已知条件，无法确定月球表面的重力加速度与地球表面重力加速度的关系，故 C 错误；由于不知道月球与地球的质量关系以及月球与地球的半径，无法比较苹果在月球表面受到的引力与在地球表面受到的引力的关系，故 D 错误。
}

\item
%% number: 18
%% paperName: 2018年北京理综第 18 题 6 分
%% typeId: 1
%% type: 单选题
%% chapter: 磁场
%% point: 带电粒子在复合场中的运动
%% method: 无
%% score: 6
%% degree: 650
%% duplicateId: 0
%% body:
某空间存在匀强磁场和匀强电场。一个带电粒子（不计重力）以一定初速度射入该空间后，做匀速直线运动；若仅撤除电场，则该粒子做匀速圆周运动，下列因素与完成上述两类运动\textbf{无关}的是\xzanswer{C}

\fourchoices[answer=C]
{磁场和电场的方向}
{磁场和电场的强弱}
{粒子的电性和电量}
{粒子入射时的速度}

%% answer: C
%% memo:
\memoanswer{
当带电粒子在复合场内做匀速直线运动时，由于不考虑粒子的重力，所以有 $qvB = qE$，即 $v = \frac{E}{B}$；当仅撤除电场，则该粒子做匀速圆周运动，洛伦兹力提供向心力。由于粒子做匀速直线运动时，粒子的电场力与洛伦兹力必须等值反向，所以与完成上述两类运动无关的是粒子的电性和电量，故 C 正确，其余项错误。
}

\item
%% number: 19
%% paperName: 2018年北京理综第 19 题 6 分
%% typeId: 1
%% type: 单选题
%% chapter: 电路
%% point: 电容
%% method: 无
%% score: 6
%% degree: 550
%% duplicateId: 0
%% body:
研究与平行板电容器电容有关因素的实验装置如图所示，下列说法正确的是\xzanswer{A}

\onepicture[w=5.5cm]{\includesvg{figs/19}}

\fourchoices[answer=A]
{实验前，只用带电玻璃棒与电容器 a 板接触，能使电容器带电}
{实验中，只将电容器 b 板向上平移，静电计指针的张角变小}
{实验中，只在极板间插入有机玻璃板，静电计指针的张角变大}
{实验中，只增加极板带电量，静电计指针的张角变大，表明电容增大}

%% answer: A
%% memo:
\memoanswer{
当用带电玻璃棒与电容器 a 板接触时，由于静电感应，从而在 b 板感应出等量的异种电荷，从而使电容器带电，故 A 正确；由平行板电容器的决定式 $C = \frac{\varepsilon S}{4\pi kd}$，将电容器 b 板向上平移，即正对面积 $S$ 减小，则电容 $C$ 减小，由 $U = \frac{Q}{C}$ 可知，电量 $Q$ 不变，则电压 $U$ 增大，则静电计指针的张角变大，故 B 错误；同理，只在极板间插入有机玻璃板，则介电系数增大，电容 $C$ 增大，电量不变时，则电压 $U$ 减小，静电计指针的张角减小，故 C 错误；电容器的电容与电量无关，故 D 错误。
}

\item
%% number: 20
%% paperName: 2018年北京理综第 20 题 6 分
%% typeId: 1
%% type: 单选题
%% chapter: 运动和力
%% point: 牛顿定律的应用
%% method: 无
%% score: 6
%% degree: 550
%% duplicateId: 0
%% body:
根据高中所学知识可知，做自由落体运动的小球，将落在正下方位置。但实际上，赤道上方 $200\Um$ 处无初速下落的小球将落在正下方位置偏东约 $6\Ucm$ 处，这一现象可解释为，除重力外，由于地球自转，下落过程小球还受到一个水平向东的“力”，该“力”与竖直方向的速度大小成正比，现将小球从赤道地面竖直上抛，考虑对称性，上升过程该“力”水平向西，则小球\xzanswer{D}

\fourchoices[answer=D]
{到最高点时，水平方向的加速度和速度均为零}
{到最高点时，水平方向的加速度和速度均不为零}
{落地点在抛出点东侧}
{落地点在抛出点西侧}

%% answer: D
%% memo:
\memoanswer{
上升过程水平方向向西加速，在最高点竖直方向上速度为零，水平方向上有向西的水平速度，且有竖直向下的加速度，故 A 错误；下降过程向西减速，按照对称性，落至地面时水平速度为 0，整个过程都在向西运动，所以落点在抛出点的西侧，故 C 错误 D 正确。
}

\item
%% number: 21
%% paperName: 2018年北京理综第 21 题 18 分
%% typeId: 4
%% type: 实验题
%% chapter: 直线运动
%% point: DIS测位移平均速度加速度
%% method: 无
%% score: 18
%% degree: 650
%% duplicateId: 0
%% body:
用图 \ref{2018北京理综21a} 所示的实验装置研究小车速度随时间变化的规律。

主要实验步骤如下：

a．安装好实验器材。接通电源后，让拖着纸带的小车沿长木板运动，重复几次。

b．选出一条点迹清晰的纸带，找一个合适的点当作计时起点 O（$t = 0$），然后每隔相同的时间间隔 $T$ 选取一个计数点，如图 \ref{2018北京理综21b} 中 A、B、C、D、E、F……所示。

c．通过测量、计算可以得到在打 A、B、C、D、E……点时小车的速度，分别记作 $v_{1}$、$v_{2}$、$v_{3}$、$v_{4}$、$v_{5}$……

d．以速度 $v$ 为纵轴、时间 $t$ 为横轴建立直角坐标系，在坐标纸上描点，如图 \ref{2018北京理综21c} 所示。

结合上述实验步骤，请你完成下列任务：
\begin{enumerate}
	\item
	在下列仪器和器材中，还需要使用的有\tkanswer[0.8cm]{A}和\tkanswer[0.8cm]{C}(填选项前的字母)。
	\item
	在图 \ref{2018北京理综21c} 中已标出计数点 A、B、D、E 对应的坐标点，请在该图中标出计数点 C 对应的坐标点，并画出 $v-t$ 图像。
	\item
	观察 $v-t$ 图像，可以判断小车做匀变速直线运动，其依据是\tkanswer[3cm]{}。$v-t$ 图像斜率的物理意义是\tkanswer[2cm]{}。
	\item
	描绘 $v-t$ 图像前，还不知道小车是否做匀变速直线运动。用平均速度 $\frac{\Delta x}{\Delta t}$ 表示各计数点的瞬时速度，从理论上讲，对 $\Delta t$ 的要求是\tkanswer[1.5cm]{}(选填“越小越好”或“与大小无关”)；从实验的角度看，选取的 $\Delta x$ 大小与速度测量的误差\tkanswer[1.5cm]{}(选填“有关”或“无关”)。
	\item
	早在 16 世纪末，伽利略就猜想落体运动的速度应该是均匀变化的。当时只能靠滴水计时，为此他设计了如图 \ref{2018北京理综21d} 所示的“斜面实验”，反复做了上百次，验证了他的猜想。请你结合匀变速直线运动的知识，分析说明如何利用伽利略“斜面实验”检验小球的速度是随时间均匀变化的\tkanswer[4cm]{}。
\end{enumerate}
\fourpicture[labelA=2018北京理综21a, labelB=2018北京理综21b, labelC=2018北京理综21c, labelD=2018北京理综21d, numA=1, numB=2, numC=3, numD=4, wA=6.5cm, wB=13cm, wC=5.5cm, wD=6.5cm, align=right]
{\includesvg{figs/21a}}
{\includesvg{figs/21b}}
{\includesvg{figs/21c}}
{\includesvg{figs/21d}}
%% answer: 
\jdanswer{
\begin{enumerate}
	\item
	AC

	\item
	如图所示：

	\onepicture[w=5.5cm]{\includesvg{figs/21_an}}

	\item
	小车的速度随时间均匀变化，加速度

	\item
	越小越好，有关

	\item
	如果小球的初速度为 0，其速度 $v \propto t$，那么它通过的位移 $x \propto t^{2}$。因此，只要测量小球通过不同位移所用的时间，就可以检验小球的速度是否随时间均匀变化。

\end{enumerate}
}
%% memo:
\memoanswer{
（1）打点计时器需用交流电源，计算速度需要利用刻度尺测量长度，故需要的仪器选 AC。

（2）利用所给点迹描点连线得到图象，其中 C 点的横坐标为 $3T$，纵坐标为 $v_{3}$。

（3）由图可知速度随时间均匀增大，所以小车做匀加速运动，图象的斜率代表了匀加速运动的加速度。

（4）当 $\Delta t$ 越小时，平均速度 $\frac{\Delta x}{\Delta t}$ 越接近各计数点的瞬时速度；由于计算速度需要用到 $\Delta x$ 的测量值，所以 $\Delta x$ 大小与速度测量的误差有关。

（5）如果小球的初速度为 0，且速度与时间成正比，那么它通过的位移 $x \propto t^{2}$。因此，只要测量小球通过不同位移所用的时间，分析 $x \propto t^{2}$ 是否成立，就可以检验小球的速度是否随时间均匀变化。
}

\item
%% number: 22
%% paperName: 2018年北京理综第 22 题 16 分
%% typeId: 6
%% type: 计算题
%% chapter: 曲线运动
%% point: 竖直平面圆周运动
%% method: 无
%% score: 16
%% degree: 750
%% duplicateId: 0
%% body:
2022 年将在我国举办第二十四届冬奥会，跳台滑雪是其中最具观赏性的项目之一。某滑道示意图如下，长直助滑道 AB 与弯曲滑道 BC 平滑衔接，滑道 BC 高 $h = 10\Um$，C 是半径 $R = 20\Um$ 圆弧的最低点，质量 $m = 60\Ukg$ 的运动员从 A 处由静止开始匀加速下滑，加速度 $a = 4.5\Umsq$，到达 B 点时速度 $v_{B} = 30\Ums$。取重力加速度 $g = 10\Umsq$。
\begin{enumerate}
	\item
	求长直助滑道 AB 的长度 $L$；
	\item
	求运动员在 AB 段所受合外力的冲量 $I$ 的大小；
	\item
	若不计 BC 段的阻力，画出运动员经过 C 点时的受力图，并求其所受支持力 $F_{N}$ 的大小。
\end{enumerate}
\onepicture[align=right, w=6.5cm]{\includesvg{figs/22}}
%% answer: 
\jdanswer{
\begin{enumerate}
	\item
	$L = 100\Um$

	\item
	$I = 1800\UNs$

	\item
	$F_{N} = 3900\UN$

\end{enumerate}
}
%% memo:
\memoanswer{
（1）由匀变速直线运动的速度与位移关系 $v^{2} - v_{0}^{2} = 2aL$，代入数据可解得 $L = 100\Um$。

（2）由动量定理可知合外力的冲量等于动量变化量，得 $I_{\text{合}} = mv_{B}$，得 $I_{\text{合}} = 1800\UNs$。

（3）运动员在最低点的受力如图所示，

\onepicture[h=3cm]{figs/22_an.png}

由 B 到 C，由动能定理得
$\frac{1}{2}mv_{C}^{2} - \frac{1}{2}mv_{B}^{2} = mgh$，
在 C 点，由牛顿第二定律得
$F_{N} - mg = \frac{mv_{C}^{2}}{R}$，
联立求得 $F_{N} = 3900\UN$。
}

\item
%% number: 23
%% paperName: 2018年北京理综第 23 题 18 分
%% typeId: 6
%% type: 计算题
%% chapter: 电路
%% point: 输出功率极值
%% method: 无
%% score: 18
%% degree: 550
%% duplicateId: 0
%% body:
如图 \ref{2018北京理综23a} 所示，用电动势为 $E$、内阻为 $r$ 的电源，向滑动变阻器 $R$ 供电。改变变阻器 $R$ 的阻值，路端电压 $U$ 与电流 $I$ 均随之变化。
\begin{enumerate}
	\item
	以 $U$ 为纵坐标，$I$ 为横坐标，在图 \ref{2018北京理综23b} 中画出变阻器阻值 $R$ 变化过程中 $U-I$ 图像的示意图，并说明 $U-I$ 图像与两坐标轴交点的物理意义。
	\item
	a．请在图 \ref{2018北京理综23b} 画好的 $U-I$ 关系图线上任取一点，画出带网格的图形，以其面积表示此时电源的输出功率；

	b．请推导该电源对外电路能够输出的最大电功率及条件。
	\item
	请写出电源电动势定义式，并结合能量守恒定律证明：电源电动势在数值上等于内、外电路电势降落之和。
\end{enumerate}
\twopicture[labelA=2018北京理综23a, labelB=2018北京理综23b, numA=1, numB=2, wA=7.5cm, wB=4cm, align=right]
{figs/23a.png}
{figs/23b.png}
%% answer: 
\jdanswer{
\begin{enumerate}
	\item
	$U-I$ 图象如图所示：

	\onepicture[w=4.5cm]{figs/23_an1.png}

	图象与纵轴交点的坐标值为电源电动势，与横轴交点的坐标值为短路电流。

	\item
	a．如图所示：

	\onepicture[w=4.5cm]{figs/23_an2.png}

	b．当 $R = r$ 时，输出功率最大，最大为 $P_{\max} = \frac{E^{2}}{4r}$。

	\item
	$E = \frac{W_{\text{非静电力}}}{q}$；由能量守恒定律可得 $E = Ir + IR = U_{\text{内}} + U_{\text{外}}$。

\end{enumerate}
}
%% memo:
\memoanswer{
（1）在图 2 中描点连线得到 $U-I$ 图象，图象与纵轴交点的坐标值为电源电动势，与横轴交点的坐标值为短路电流。

（2）a．在 $U-I$ 关系图线上任取一点，画出带网格的图形，以其面积表示此时电源的输出功率。

b．由闭合电路欧姆定律得 $I = \frac{E}{R + r}$，输出功率 $P = I^{2}R = \frac{E^{2}}{R + \frac{r^{2}}{R} + 2r}$，由于 $R + \frac{r^{2}}{R} \geqslant 2r$，所以当 $R = r$ 时，输出功率最大，最大为 $P_{\max} = \frac{E^{2}}{4r}$。

（3）电动势的定义式 $E = \frac{W_{\text{非静电力}}}{q}$，由能量守恒定律，非静电力做的功等于内、外电路上产生的电热，有 $W = I^{2}rt + I^{2}Rt = Irq + IRq$，即 $IEt = I^{2}Rt + I^{2}rt$，得到 $E = Ir + IR = U_{\text{内}} + U_{\text{外}}$。
}

\item
%% number: 24
%% paperName: 2018年北京理综第 24 题 20 分
%% typeId: 6
%% type: 计算题
%% chapter: 电场
%% point: 电场强度
%% method: 建立模型
%% score: 20
%% degree: 450
%% duplicateId: 0
%% body:
\begin{enumerate}
	\item
	静电场可以用电场线和等势面形象描述。

	a．请根据电场强度的定义和库仑定律推导出点电荷 $Q$ 的场强表达式；

	b．点电荷的电场线和等势面分布如图所示，等势面 $S_{1}$、$S_{2}$ 到点电荷的距离分别为 $r_{1}$、$r_{2}$。我们知道，电场线的疏密反映了空间区域电场强度的大小。请计算 $S_{1}$、$S_{2}$ 上单位面积通过的电场线条数之比 $N_{1}/N_{2}$。
	\item
	观测宇宙中辐射电磁波的天体，距离越远单位面积接收的电磁波功率越小，观测越困难。为了收集足够强的来自天体的电磁波，增大望远镜口径是提高天文观测能力的一条重要路径。2016 年 9 月 25 日，世界上最大的单口径球面射电望远镜 FAST 在我国贵州落成启用，被誉为“中国天眼”。FAST 直径为 $500\Um$，有效提高了人类观测宇宙的精度和范围。

	a．设直径为 $100\Um$ 的望远镜能够接收到的来自某天体的电磁波功率为 $P_{1}$，计算 FAST 能够接收到的来自该天体的电磁波功率 $P_{2}$；

	b．在宇宙大尺度上，天体的空间分布是均匀的，仅以辐射功率为 $P$ 的同类天体为观测对象，设直径为 $100\Um$ 望远镜能够观测到的此类天体数目是 $N_{0}$，计算 FAST 能够观测到的此类天体数目 $N$。
\end{enumerate}
\onepicture[align=right, w=5.5cm]{\includesvg{figs/24}}
%% answer: 
\jdanswer{
\begin{enumerate}
	\item
	a．设检验电荷的电量为 $q$，距离点电荷 $Q$ 为 $r$，受到点电荷 $Q$ 的作用力为 $F$，由 $E = \frac{F}{q}$ 以及库仑定律 $F = \frac{kQq}{r^{2}}$，可得 $E = \frac{kQ}{r^{2}}$。

	b．$\frac{N_{1}}{N_{2}} = \frac{E_{1}}{E_{2}} = \frac{r_{2}^{2}}{r_{1}^{2}}$
	\item
	a．$P_{2} = \frac{500^{2}}{100^{2}}P_{1} = 25P_{1}$

	b．$N = 125N_{0}$

\end{enumerate}
}
%% memo:
\memoanswer{
（1）a．由库仑定律以及电场强度的定义式，推导点电荷的场强公式。设检验电荷的电量为 $q$，距离点电荷 $Q$ 为 $r$，受到的作用力为 $F$，由 $E = \frac{F}{q}$ 及 $F = \frac{kQq}{r^{2}}$，可得 $E = \frac{kQ}{r^{2}}$。

b．由穿过两等势面单位面积上的电场线条数之比等于电场强度之比，可求得
$\frac{N_{1}}{N_{2}} = \frac{E_{1}}{E_{2}} = \frac{r_{2}^{2}}{r_{1}^{2}}$。

（2）a．地球上不同望远镜观测同一天体，单位面积上接收的功率相同，所以
$P_{2} = \frac{500^{2}}{100^{2}}P_{1} = 25P_{1}$。

b．在宇宙大尺度上，天体的空间分布是均匀的，因此，一个望远镜能观测的此类天体数目正比于以望远镜为球心、以最远观测距离为半径的球体的体积。设地面上望远镜能观测到此类天体需收集到的电磁波的总功率的最小值为 $P_{0}$，直径为 $100\Um$ 望远镜和 FAST 能观测到的最远距离分别为 $L_{0}$ 和 $L$，则
$P_{0} = \pi\left(\frac{500}{2}\right)^{2}\frac{P}{4\pi L^{2}} = \pi\left(\frac{100}{2}\right)^{2}\frac{P}{4\pi L_{0}^{2}}$，
可得 $L = 5L_{0}$，则
$N = \frac{L^{3}}{L_{0}^{3}}N_{0} = 125N_{0}$。
}

\end{enumerate}

\end{document}
